\section{Results}\label{sec:results}

\subsection{H1: FFL Motif Universality}\label{sec:results_h1}

Across all 200 attribution graphs and 50 degree-preserving null models per graph, the FFL motif (030T) achieves a median Z-score of $47.2$ (mean $Z=46.2$, $\sigma=12.0$), with 100\% of graphs exceeding $Z>2$ (Table~\ref{table:zscores}). This represents an extremely strong signal: the observed FFL count exceeds the null expectation by tens of standard deviations in every single graph. All 8 capability domains show positive FFL over-representation, with per-domain median Z-scores ranging from 26.3 (antonym) to 59.4 (sentiment). Notably, even the weakest domain (antonym) shows a Z-score far exceeding conventional significance thresholds. The three non-FFL 3-node DAG motifs (021U fan-out, 021C chain, 021D fan-in) are uniformly \emph{under}-represented, with negative Z-scores of comparable magnitude. This asymmetry --- FFL over-representation coupled with non-FFL under-representation --- is consistent with the hypothesis that attribution circuits preferentially organize computation through feedforward integration rather than simple fan-out or fan-in patterns.

\paragraph{FDR Correction and Honest Assessment.}
Benjamini--Hochberg correction at $\alpha=0.05$ across 800 3-node tests yields 0/800 rejections post-FDR, and 0/3600 for 4-node tests. This apparent contradiction --- extreme Z-scores yet zero FDR survival --- is an important methodological finding. It arises because the limited number of null models ($N=50$) produces discrete p-value distributions: the minimum attainable p-value per test is $1/50=0.020$, which does not survive strict multiple-testing correction across hundreds of tests. This is a limitation of the null model sample size, not evidence against universality. The Z-score magnitudes themselves are unambiguous and would require astronomical numbers of null models to produce survival under BH-FDR. We present this transparently as a methodological caveat.



\paragraph{Corpus-Level Statistical Tests.}
Aggregating Z-scores across all 200 graphs, a one-sample $t$-test yields $t = 55.03$ ($p = 2.38 \times 10^{-122}$), with Cohen's $d = 3.89$ indicating an extremely large effect size. The Wilcoxon signed-rank test confirms significance ($p = 1.44 \times 10^{-34}$), and a sign test finds all 200/200 graphs show positive FFL Z-scores ($p < 10^{-60}$). A mixed-effects model with domain as random intercept estimates $\beta_0 = 47.18$ (95\% CI: [40.54, 53.82]), with ICC $= 0.570$ indicating that 57\% of Z-score variance is between-domain. Per-graph BH-FDR with 200 null models on 15 stratified graphs yields 15/60 tests surviving correction, confirming significance at the individual-graph level when null resolution is adequate. Comparing to biological networks, LLM FFL Z-scores ($\bar{Z} = 47.14$) are 3.7--5.5$\times$ larger than canonical biological FFLs (\emph{E.~coli} $Z \approx 12.7$, yeast $Z \approx 8.5$).
\paragraph{Convergence and Robustness.}
Convergence diagnostics (Phase~E) confirm that Z-score estimates stabilize at a median of $N=30$ null models, well below our budget of 50. By $N=50$, 100\% of tested graphs have converged to within 10\% of their final Z-score estimates. Pruning robustness analysis (Phase~D) across three thresholds (60th, 75th, 90th percentiles) confirms stable motif rankings despite varying graph density. Layer-preserving null models (Phase~F), which rewire edges while maintaining the DAG layer structure, yield even stronger FFL Z-scores (mean $Z=56.9$, an increase of $10.7$ relative to standard degree-preserving nulls, $p<0.001$ by Wilcoxon signed-rank test). This is a crucial robustness check: it demonstrates that FFL over-representation is not an artifact of the layered DAG structure inherent to attribution graphs, but reflects genuine local circuit organization.

These findings are consistent with the motif universality hypothesis of \citet{milo2002network}, who observed that the FFL is over-represented across diverse biological and technological networks. Our results extend this finding to artificial neural network circuits, suggesting that the FFL captures a fundamental information-processing primitive that emerges independently in both biological evolution and gradient-based optimization \cite{alon2007network, milo2004superfamilies}.

\textbf{Verdict:} H1 is \textbf{supported}. FFL motifs are universally over-represented across all 8/8 capability domains with large effect sizes, robust to pruning threshold, null model choice, and layer structure controls.